%% Drafted from the mapped corpus by scripts/draft_topics.py.
% Model: gpt-5.6-terra; source characters: 20622.
\section{Симетрична група и цикличен запис}

\begin{defn}
Нека $M$ е множество. Множеството
\[
\operatorname{Sym}(M)=\{f\colon M\to M\mid f\text{ е биекция}\}
\]
е група относно композицията на изображенията. Тя се нарича \emph{симетрична група на множеството $M$}. Неутралният ѝ елемент е тъждественото изображение
\[
\operatorname{Id}_M\colon M\to M,\qquad \operatorname{Id}_M(x)=x.
\]
\end{defn}

Ако $|M|=n$, то след избиране на номерация $M=\{1,\ldots,n\}$ симетричната група се бележи със $S_n$ и се нарича \emph{симетрична група от степен $n$}. Нейните елементи се наричат \emph{пермутации}. Всяка пермутация е биекция на $\{1,\ldots,n\}$, следователно се определя еднозначно от редицата
\[
\sigma(1),\sigma(2),\ldots,\sigma(n),
\]
която е пренареждане на числата $1,\ldots,n$. Поради това
\[
|S_n|=n!.
\]

\begin{defn}
За $\sigma\in S_n$ множеството
\[
\operatorname{Supp}(\sigma)=\{i\in\{1,\ldots,n\}\mid \sigma(i)\ne i\}
\]
се нарича \emph{носител} на $\sigma$.
\end{defn}

\begin{defn}
Пермутацията $\xi\in S_n$ е \emph{цикъл с дължина $k$}, ако има $k$ различни числа
\[
i_1,\ldots,i_k
\]
такива, че
\[
\xi(i_s)=i_{s+1}\quad (1\le s<k),\qquad \xi(i_k)=i_1,
\]
а всички останали числа се оставят неподвижни. Тя се записва като
\[
\xi=(i_1,i_2,\ldots,i_k).
\]
\end{defn}

Един и същ цикъл има много записи, получени чрез циклично преместване:
\[
(i_1,i_2,\ldots,i_k)
=
(i_2,\ldots,i_k,i_1).
\]
Редът на елементите в цикличния запис обаче е съществен: обикновено
\[
(i_1,i_2,\ldots,i_k)\ne(i_k,\ldots,i_2,i_1).
\]

\begin{defn}
Два цикъла се наричат \emph{независими}, ако носителите им не се пресичат.
\end{defn}

Независимите цикли комутират, тъй като всеки от тях оставя неподвижни всички елементи от носителя на другия. Следователно редът на множителите в произведение на независими цикли няма значение.

\begin{thm}
Всяка нетъждествена пермутация $\sigma\in S_n\setminus\{\operatorname{Id}\}$ се представя като произведение на независими цикли. Представянето е единствено с точност до реда на множителите и до цикличното преместване в записа на всеки отделен цикъл.
\end{thm}

\begin{proof}
Избираме $i_1\in\operatorname{Supp}(\sigma)$ и разглеждаме последователността
\[
i_1,\ \sigma(i_1),\ \sigma^2(i_1),\ldots.
\]
Тъй като множеството $\{1,\ldots,n\}$ е крайно, някой член трябва да се повтори. Първото повторение е именно $i_1$: ако $\sigma(i_r)=i_s$ за $1\le s<r$, то от инективността на $\sigma$ следва $i_{r-1}=i_{s-1}$, което противоречи на избора на първо повторение, освен когато $s=1$. Получава се цикъл
\[
(i_1,\sigma(i_1),\ldots,\sigma^{k-1}(i_1)).
\]

Множеството от елементи на този цикъл се запазва от $\sigma$, а върху допълнението му $\sigma$ отново е пермутация. Повтаряйки процедурата върху останалите неподвижни още елементи от носителя, получаваме разлагане на $\sigma$ в независими цикли.

За единствеността нека имаме две разлагания в независими цикли. Ако $i$ принадлежи на цикъл от първото разлагане, то в другото разлагане има точно един цикъл, който мести $i$. И в двете разлагания последователните образи
\[
i,\ \sigma(i),\ \sigma^2(i),\ldots
\]
са едни и същи; следователно двата цикъла съвпадат. След премахването им разсъждението се прилага за останалите цикли.
\end{proof}

\begin{example}
Нека $\sigma\in S_8$ е зададена чрез
\[
\sigma(1)=4,\quad \sigma(4)=7,\quad \sigma(7)=1,
\]
\[
\sigma(2)=6,\quad \sigma(6)=2,\quad \sigma(3)=3,\quad \sigma(5)=5,\quad \sigma(8)=8.
\]
Тогава
\[
\sigma=(1,4,7)(2,6).
\]
Фиксираните точки $3$, $5$ и $8$ обикновено не се записват.
\end{example}

\section{Спрягане в $S_n$}

\begin{defn}
Елементите $a$ и $b$ на група $G$ се наричат \emph{спрегнати}, ако съществува $c\in G$, за което
\[
b=cac^{-1}.
\]
\end{defn}

В симетричната група спрягането има особено ясен смисъл: то просто преименува елементите, участващи в цикъла.

\begin{lem}
За $\rho\in S_n$ и цикъл $(i_1,\ldots,i_k)\in S_n$ е изпълнено
\[
\rho(i_1,\ldots,i_k)\rho^{-1}
=
\bigl(\rho(i_1),\ldots,\rho(i_k)\bigr).
\]
\end{lem}

\begin{proof}
Нека $1\le s\le k$, като индексите в цикъла се разглеждат по модул $k$. Тогава
\[
\bigl[\rho(i_1,\ldots,i_k)\rho^{-1}\bigr](\rho(i_s))
=
\rho(i_1,\ldots,i_k)(i_s)
=
\rho(i_{s+1}).
\]
Това е точно действието на цикъла
\[
\bigl(\rho(i_1),\ldots,\rho(i_k)\bigr).
\]
Ако $j$ не участва в първоначалния цикъл, то $\rho(j)$ не участва в получения цикъл и и двете страни оставят $\rho(j)$ неподвижно.
\end{proof}

Ако
\[
\sigma=\alpha_1\alpha_2\cdots\alpha_r
\]
е разлагане в независими цикли, тогава
\[
\rho\sigma\rho^{-1}
=
(\rho\alpha_1\rho^{-1})
(\rho\alpha_2\rho^{-1})
\cdots
(\rho\alpha_r\rho^{-1}).
\]
Следователно при спрягане се запазват дължините и броят на независимите цикли.

\begin{defn}
Две пермутации от $S_n$ имат \emph{еднакъв цикличен строеж}, ако в разлаганията си в независими цикли имат еднакъв брой цикли от всяка дължина.
\end{defn}

\begin{prop}
Две пермутации $\sigma,\tau\in S_n$ са спрегнати в $S_n$ тогава и само тогава, когато имат еднакъв цикличен строеж.
\end{prop}

\begin{proof}
Ако $\tau=\rho\sigma\rho^{-1}$, лемата показва, че всеки цикъл на $\sigma$ се преобразува в цикъл със същата дължина. Значи цикличният строеж се запазва.

Обратно, нека $\sigma$ и $\tau$ имат еднакъв цикличен строеж. Сдвояваме циклите им с еднаква дължина и дефинираме $\rho$ така, че да изпраща последователните елементи на всеки цикъл на $\sigma$ в последователните елементи на съответния цикъл на $\tau$. Върху фиксираните точки $\rho$ се избира произволно като биекция към фиксираните точки на $\tau$. Тогава лемата дава
\[
\tau=\rho\sigma\rho^{-1}.
\]
\end{proof}

\section{Транспозиции, четност и алтернативна група}

\begin{defn}
Цикъл с дължина $2$ се нарича \emph{транспозиция}. Той има вида
\[
(i,j)
\]
и разменя $i$ и $j$, като оставя всички останали елементи неподвижни.
\end{defn}

За всяка транспозиция е изпълнено
\[
(i,j)^{-1}=(i,j),\qquad (i,j)^2=\operatorname{Id}.
\]

\begin{prop}
Всеки цикъл с дължина $k\ge2$ се представя като произведение на $k-1$ транспозиции:
\[
(i_1,i_2,\ldots,i_k)
=
(i_1,i_k)(i_1,i_{k-1})\cdots(i_1,i_3)(i_1,i_2).
\]
\end{prop}

\begin{proof}
Прилагаме множителите отдясно наляво. Елементът $i_1$ се изпраща в $i_2$ от първата транспозиция и следващите го фиксират. За $2\le s<k$ елементът $i_s$ остава неподвижен до $(i_1,i_s)$, която го изпраща в $i_1$; следващата $(i_1,i_{s+1})$ го изпраща в $i_{s+1}$, където остава. Елементът $i_k$ се изпраща в $i_1$ от последната транспозиция. Всички останали елементи са фиксирани. Това е точно действието на цикъла.
\end{proof}



\begin{cor}
Всяка пермутация от $S_n$ се представя като произведение на транспозиции.
\end{cor}

\begin{proof}
Разлагаме пермутацията на независими цикли по доказаната теорема. Заменяме всеки цикъл с дължина поне $2$ с произведението от транспозиции от предходното твърдение. Неподвижните точки не изискват множител; тъждествената пермутация се представя с празното произведение.
\end{proof}

Наистина, първо разлагаме пермутацията на независими цикли, след което всеки цикъл заменяме с произведението от транспозиции, дадено по-горе.

\begin{thm}
Ако
\[
\tau_1\tau_2\cdots\tau_k=\operatorname{Id},
\]
където $\tau_1,\ldots,\tau_k$ са транспозиции, то $k$ е четно число. Следователно ако една и съща пермутация е представена като произведение съответно на $k$ и на $m$ транспозиции, то
\[
k\equiv m\pmod 2.
\]
\end{thm}

\begin{proof}
На пермутация $\sigma$ съпоставяме пермутационната матрица $P_\sigma$, която изпраща $e_i$ в $e_{\sigma(i)}$. Тогава $P_{\sigma\rho}=P_\sigma P_\rho$. Матрицата на транспозиция се получава от единичната с размяна на две колони, затова детерминантата й е $-1$. От произведението, равно на тъждествената пермутация, получаваме $1=(-1)^k$, следователно $k$ е четно. Ако два записа на една пермутация имат дължини $k,m$, умножаваме първия по обратния на втория. Получаваме запис на тъждествената с $k+m$ транспозиции, така че $k+m$ е четно, тоест $k\equiv m\pmod2$.
\end{proof}

Теоремата позволява да се въведе четността на пермутация независимо от избраното разлагане на транспозиции.

\begin{defn}
Пермутацията $\sigma\in S_n$ се нарича \emph{четна}, ако в произволно нейно представяне като произведение на транспозиции броят на транспозициите е четен. Тя се нарича \emph{нечетна}, ако този брой е нечетен.
\end{defn}

Тъждествената пермутация е четна. Произведението на две четни пермутации е четно; обратната на четна пермутация също е четна. Аналогично, произведението на четна и нечетна пермутация е нечетно, а произведението на две нечетни пермутации е четно.

\begin{remark}
Нека $\sigma\in S_n$ и разгледаме редицата
\[
\sigma(1),\ldots,\sigma(n).
\]
Двойката $\sigma(p),\sigma(q)$ образува \emph{инверсия}, ако
\[
1\le p<q\le n
\quad\text{и}\quad
\sigma(p)>\sigma(q).
\]
Четността на пермутацията съвпада с четността на броя на нейните инверсии.
\end{remark}

\begin{defn}
За $n\ge2$ множеството
\[
A_n=\{\sigma\in S_n\mid \sigma\text{ е четна}\}
\]
се нарича \emph{алтернативна група от степен $n$}.
\end{defn}

\begin{thm}
За всяко $n\ge2$ множеството $A_n$ е нормална подгрупа на $S_n$ с индекс $2$. В частност
\[
A_n\triangleleft S_n,
\qquad
[S_n:A_n]=2,
\qquad
|A_n|=\frac{n!}{2}.
\]
\end{thm}

\begin{proof}
Критерият за подгрупа показва, че $A_n\le S_n$: ако $\sigma,\tau\in A_n$, то $\sigma\tau^{-1}$ е произведение на четен брой транспозиции и следователно принадлежи на $A_n$.

Нека $(i,j)$ е фиксирана транспозиция. Произведението $(i,j)\sigma$ е нечетно за всяка $\sigma\in A_n$, така че
\[
(i,j)A_n\subseteq S_n\setminus A_n.
\]
Обратно, ако $\tau$ е нечетна, то $(i,j)\tau$ е четна. Следователно
\[
\tau=(i,j)\bigl((i,j)\tau\bigr)\in(i,j)A_n.
\]
Така
\[
S_n\setminus A_n=(i,j)A_n.
\]
Получаваме разлагането на $S_n$ в два леви съседни класа:
\[
S_n=A_n\mathbin{\dot\cup}(i,j)A_n.
\]
Следователно индексът е $2$. Подгрупа с индекс $2$ е нормална, а от теоремата на Лагранж следва
\[
|A_n|=\frac{|S_n|}{[S_n:A_n]}=\frac{n!}{2}.
\]
\end{proof}

\section{Хомоморфизми, ядро и образ}

\begin{defn}
Нека $(G_1,\cdot)$ и $(G_2,\ast)$ са групи. Отображението
\[
\varphi\colon G_1\to G_2
\]
се нарича \emph{хомоморфизъм на групи}, ако за всички $a,b\in G_1$ е изпълнено
\[
\varphi(a\cdot b)=\varphi(a)\ast\varphi(b).
\]
\end{defn}

Ако груповите операции са означени еднакво, обичайно се пише просто
\[
\varphi(ab)=\varphi(a)\varphi(b).
\]

\begin{prop}
Нека $\varphi\colon G_1\to G_2$ е хомоморфизъм, а $e_1$ и $e_2$ са неутралните елементи на $G_1$ и $G_2$. Тогава:
\[
\varphi(e_1)=e_2;
\]
\[
\varphi(a^{-1})=\varphi(a)^{-1}
\qquad\text{за всяко }a\in G_1.
\]
\end{prop}

\begin{proof}
От $e_1e_1=e_1$ получаваме
\[
\varphi(e_1)=\varphi(e_1)\varphi(e_1).
\]
Умножаваме отляво с $\varphi(e_1)^{-1}$ и получаваме $\varphi(e_1)=e_2$.

Освен това
\[
e_2=\varphi(e_1)=\varphi(aa^{-1})
=\varphi(a)\varphi(a^{-1}),
\]
поради което $\varphi(a^{-1})$ е обратен на $\varphi(a)$.
\end{proof}

\begin{defn}
За хомоморфизъм $\varphi\colon G_1\to G_2$ множествата
\[
\ker\varphi
=
\{a\in G_1\mid \varphi(a)=e_2\}
\]
и
\[
\operatorname{Im}\varphi
=
\{\varphi(a)\mid a\in G_1\}
\]
се наричат съответно \emph{ядро} и \emph{образ} на $\varphi$.
\end{defn}

\begin{prop}
Ако $\varphi\colon G_1\to G_2$ е хомоморфизъм, то
\[
\ker\varphi\le G_1,
\qquad
\operatorname{Im}\varphi\le G_2.
\]
\end{prop}

\begin{proof}
За $a,b\in\ker\varphi$ имаме
\[
\varphi(ab^{-1})
=
\varphi(a)\varphi(b)^{-1}
=
e_2e_2^{-1}
=
e_2.
\]
Следователно $ab^{-1}\in\ker\varphi$ и критерият за подгрупа доказва първото твърдение.

Ако $x=\varphi(a)$ и $y=\varphi(b)$ принадлежат на $\operatorname{Im}\varphi$, то
\[
xy^{-1}
=
\varphi(a)\varphi(b)^{-1}
=
\varphi(ab^{-1})\in\operatorname{Im}\varphi.
\]
Следователно и образът е подгрупа.
\end{proof}

За подгрупа $H\le G$ левият съседен клас с представител $a\in G$ е
\[
aH=\{ah\mid h\in H\}.
\]
Ако $H$ е нормална подгрупа, пишем $H\triangleleft G$ и левите и десните съседни класове съвпадат. Тогава множеството
\[
G/H=\{aH\mid a\in G\}
\]
е факторгрупа относно операцията
\[
(aH)(bH)=abH.
\]

\begin{prop}
Нека $\varphi\colon G_1\to G_2$ е хомоморфизъм и $H=\ker\varphi$. Тогава:
\[
\varphi(a)=\varphi(b)
\quad\Longleftrightarrow\quad
aH=bH
\]
за всички $a,b\in G_1$.
\end{prop}

\begin{proof}
Имаме
\[
\varphi(a)=\varphi(b)
\quad\Longleftrightarrow\quad
\varphi(a^{-1}b)=e_2
\quad\Longleftrightarrow\quad
a^{-1}b\in H.
\]
Последното условие е еквивалентно на $b\in aH$, а то от своя страна е еквивалентно на $aH=bH$.
\end{proof}

\section{Теорема за хомоморфизмите}

\begin{keythm}[Теорема за хомоморфизмите на групи]
Нека
\[
\varphi\colon G_1\to G_2
\]
е хомоморфизъм. Тогава
\[
\ker\varphi\triangleleft G_1
\]
и $\varphi$ индуцира изоморфизъм
\[
\overline{\varphi}\colon G_1/\ker\varphi\longrightarrow\operatorname{Im}\varphi,
\qquad
\overline{\varphi}(a\ker\varphi)=\varphi(a).
\]
За естествения хомоморфизъм
\[
\pi\colon G_1\to G_1/\ker\varphi,
\qquad
\pi(a)=a\ker\varphi,
\]
е изпълнено
\[
\overline{\varphi}\circ\pi=\varphi.
\]
Следователно
\[
G_1/\ker\varphi\cong\operatorname{Im}\varphi.
\]
\end{keythm}

\begin{proof}
За $a\in G_1$ и $h\in\ker\varphi$ имаме
\[
\varphi(aha^{-1})
=
\varphi(a)\varphi(h)\varphi(a)^{-1}
=
\varphi(a)e_2\varphi(a)^{-1}
=
e_2.
\]
Следователно $aha^{-1}\in\ker\varphi$, тоест $\ker\varphi\triangleleft G_1$.

Дефинираме
\[
\overline{\varphi}(a\ker\varphi)=\varphi(a).
\]
Това е коректно дефинирано, защото от
\[
a\ker\varphi=b\ker\varphi
\]
следва $\varphi(a)=\varphi(b)$.

За $a,b\in G_1$:
\[
\overline{\varphi}\bigl((a\ker\varphi)(b\ker\varphi)\bigr)
=
\overline{\varphi}(ab\ker\varphi)
=
\varphi(ab)
=
\varphi(a)\varphi(b),
\]
тоест $\overline{\varphi}$ е хомоморфизъм.

Той е сюрективен по дефиницията на образа: всеки елемент на $\operatorname{Im}\varphi$ има вид $\varphi(a)$ и е образ на $a\ker\varphi$. Ако
\[
\overline{\varphi}(a\ker\varphi)
=
\overline{\varphi}(b\ker\varphi),
\]
то $\varphi(a)=\varphi(b)$, откъдето $a\ker\varphi=b\ker\varphi$. Значи $\overline{\varphi}$ е и инективен.

Накрая:
\[
(\overline{\varphi}\circ\pi)(a)
=
\overline{\varphi}(a\ker\varphi)
=
\varphi(a).
\]
\end{proof}

\begin{example}
Нека
\[
\varphi\colon \mathbb{Z}\to\mathbb{Z}_m,
\qquad
\varphi(k)=\overline{k},
\]
където $\mathbb{Z}$ е адитивната група на целите числа, а $\mathbb{Z}_m$ е адитивната група по модул $m$. Тогава
\[
\ker\varphi=m\mathbb{Z},
\qquad
\operatorname{Im}\varphi=\mathbb{Z}_m.
\]
Теоремата за хомоморфизмите дава
\[
\mathbb{Z}/m\mathbb{Z}\cong\mathbb{Z}_m.
\]
\end{example}

\section{Теорема на Кейли}

\begin{keythm}[Теорема на Кейли]
Всяка крайна група $G$ от ред $n$ е изоморфна на подгрупа на симетричната група $S_n$.
\end{keythm}

\begin{proof}
Разглеждаме действието на $G$ върху самата себе си чрез леви транслации. За всяко $a\in G$ дефинираме
\[
\lambda_a\colon G\to G,
\qquad
\lambda_a(x)=ax.
\]
Отображението $\lambda_a$ е биекция, защото
\[
\lambda_a^{-1}=\lambda_{a^{-1}}.
\]
Следователно $\lambda_a\in\operatorname{Sym}(G)$.

Нека
\[
\Phi\colon G\to\operatorname{Sym}(G),
\qquad
\Phi(a)=\lambda_a.
\]
За $a,b,x\in G$ е изпълнено
\[
(\Phi(a)\circ\Phi(b))(x)
=
\Phi(a)(bx)
=
a(bx)
=
(ab)x
=
\Phi(ab)(x).
\]
Следователно $\Phi$ е хомоморфизъм.

Неговото ядро е тривиално. Действително, ако $\Phi(a)=\operatorname{Id}_G$, то
\[
ab=b
\qquad\text{за всяко }b\in G.
\]
За $b=e$ получаваме $a=e$. Значи
\[
\ker\Phi=\{e\}.
\]
По теоремата за хомоморфизмите
\[
G/\{e\}\cong\operatorname{Im}\Phi.
\]
Но $G/\{e\}\cong G$, а $\operatorname{Im}\Phi$ е подгрупа на $\operatorname{Sym}(G)$. Тъй като $|G|=n$, имаме
\[
\operatorname{Sym}(G)\cong S_n.
\]
Следователно $G$ е изоморфна на подгрупа на $S_n$.
\end{proof}

\begin{supp}[Бележка]
По-общото твърдение, доказано чрез левите транслации, е че всяка група $G$ е изоморфна на подгрупа на $\operatorname{Sym}(G)$. Формулировката с $S_n$ е специализацията за крайна група с $n$ елемента.
\end{supp}

\section{Изпитен фокус}

\begin{itemize}
\item Да се дадат определенията за $\operatorname{Sym}(M)$, $S_n$, носител на пермутация, цикъл, независими цикли, транспозиция, четна и нечетна пермутация, $A_n$, хомоморфизъм, ядро, образ, нормална подгрупа и факторгрупа.

\item Да се обясни алгоритъмът за разлагане на пермутация в независими цикли: започва се от елемент от носителя, проследяват се последователните му образи до завръщане в началния елемент и процедурата се повтаря върху останалите елементи.

\item Да се възпроизведе формулата
\[
(i_1,\ldots,i_k)
=
(i_1,i_k)(i_1,i_{k-1})\cdots(i_1,i_2)
\]
и изводът, че всяка пермутация е произведение на транспозиции.

\item Да се знае, че броят на транспозициите в две разлагания на една и съща пермутация има една и съща четност, и да се използва това за определяне на четни и нечетни пермутации.

\item Да се формулира и докаже, че
\[
A_n\triangleleft S_n,\qquad [S_n:A_n]=2,\qquad |A_n|=\frac{n!}{2}.
\]

\item Да се използва формулата за спрягане на цикъл:
\[
\rho(i_1,\ldots,i_k)\rho^{-1}
=
(\rho(i_1),\ldots,\rho(i_k)).
\]
Да се обясни критерият: две пермутации в $S_n$ са спрегнати точно когато имат еднакъв цикличен строеж.

\item Да се докажат основните свойства на хомоморфизма:
\[
\varphi(e_1)=e_2,
\qquad
\varphi(a^{-1})=\varphi(a)^{-1},
\qquad
\ker\varphi\le G_1,
\qquad
\operatorname{Im}\varphi\le G_2.
\]

\item Да се формулира теоремата за хомоморфизмите:
\[
G_1/\ker\varphi\cong\operatorname{Im}\varphi,
\]
като се обясни защо отображението
\[
a\ker\varphi\longmapsto\varphi(a)
\]
е коректно дефинирано, хомоморфизъм, сюрективно и инективно.

\item Да се възпроизведе доказателствената идея на теоремата на Кейли: групата действа върху себе си чрез леви транслации $x\mapsto ax$; получава се хомоморфизъм към симетрична група с тривиално ядро.
\end{itemize}
